% GENERATED FILE. Edit canonical lessons, then run npm run generate:books.
% canonical-source-hash: fnv1a64:beba478294109551
% canonical-lessons: GE-C354-besprechung, GE-C354-sitzung, GE-C354-protokoll, GE-C354-praesentieren, GE-C354-praesentation

\chapter{Die Besprechung}
\label{ch:ge-die-besprechung}
\hlchaptermodality{354}

\begin{chapteropening}
\textbf{By the end of this chapter:} \emph{I can talk about meetings at work: who takes the minutes and who gives a presentation.}

Complete the last lesson of chapter 354: I can talk about meetings at work: who takes the minutes and who gives a presentation.
\end{chapteropening}

\section[die Besprechung]{die Besprechung — the meeting (at work)}

\label{lesson:GE-C354-besprechung}



\begin{quote}
Before the new one: say the German for the superior, then the German for to run.
\end{quote}

\subsection*{You'll want to know: die Besprechung}
\textbf{die Besprechung} — \textquotedblleft{}the meeting (at work)\textquotedblright{}.

\begin{quote}
Die Besprechung dauert bis elf Uhr.
\end{quote}

\begin{quote}
\emph{The meeting lasts until eleven o'clock.}
\end{quote}

Plural: \emph{die Besprechungen}. \emph{In einer Besprechung sein} — to be in a meeting.

\begin{cousinweb}
\emph{Besprechung} is the noun of \emph{besprechen} \textquotedblleft{}to discuss\textquotedblright{}, built with \emph{be-} on \emph{sprechen} \textquotedblleft{}to speak\textquotedblright{}, Old High German \emph{sprehhan} and Proto-Germanic \emph{sprekaną} (EtymDB).
\end{cousinweb}

\begin{grammarlens}[title={What you've built}]
One more word for telling, explaining and describing at B1.
\end{grammarlens}

\subsection*{Your turn}
\begin{itemize}
\raggedright
  \item \emph{Say it:} \emph{die Besprechung}
  \item \emph{Say it:} \emph{die Besprechung}, once more
  \item \emph{Say it:} a sentence of your own with \emph{die Besprechung}
  \item \emph{Recall:} say the German for the superior, then the German for to run, then say \emph{die Besprechung} again
\end{itemize}

\begin{tcolorbox}[breakable,colback=teal!4,colframe=teal!35!black,title={Before you move on}]
What does die Besprechung mean? (\textquotedblleft{}The meeting (at work)\textquotedblright{}.) Say it once more.
\end{tcolorbox}
\section[die Sitzung]{die Sitzung — the session, the formal meeting}

\label{lesson:GE-C354-sitzung}



\begin{quote}
Before the new one: say the German for to run, then the German for the meeting.
\end{quote}

\subsection*{You'll want to know: die Sitzung}
\textbf{die Sitzung} — \textquotedblleft{}the session, the formal meeting\textquotedblright{}.

\begin{quote}
Der Stadtrat hat jeden Montag eine Sitzung.
\end{quote}

\begin{quote}
\emph{The town council has a session every Monday.}
\end{quote}

Plural: \emph{die Sitzungen}. More formal than \emph{die Besprechung}; with a therapist, a session.

\begin{cousinweb}
\emph{Sitzung} is formed from \emph{sitzen} \textquotedblleft{}to sit\textquotedblright{}, inherited from Proto-Germanic \emph{sitjaną} (EtymDB), with \emph{-ung}: a sitting, as in English.
\end{cousinweb}

\begin{grammarlens}[title={What you've built}]
One more word for telling, explaining and describing at B1.
\end{grammarlens}

\subsection*{Your turn}
\begin{itemize}
\raggedright
  \item \emph{Say it:} \emph{die Sitzung}
  \item \emph{Say it:} \emph{die Sitzung}, once more
  \item \emph{Say it:} a sentence of your own with \emph{die Sitzung}
  \item \emph{Recall:} say the German for to run, then the German for the meeting, then say \emph{die Sitzung} again
\end{itemize}

\begin{tcolorbox}[breakable,colback=teal!4,colframe=teal!35!black,title={Before you move on}]
What does die Sitzung mean? (\textquotedblleft{}The session, the formal meeting\textquotedblright{}.) Say it once more.
\end{tcolorbox}
\section[das Protokoll]{das Protokoll — the minutes, the record; protocol}

\label{lesson:GE-C354-protokoll}



\begin{quote}
Before the new one: say the German for the meeting, then the German for the session.
\end{quote}

\subsection*{You'll want to know: das Protokoll}
\textbf{das Protokoll} — \textquotedblleft{}the minutes, the record; protocol\textquotedblright{}.

\begin{quote}
Heute führt mein Kollege das Protokoll.
\end{quote}

\begin{quote}
\emph{Today my colleague is taking the minutes.}
\end{quote}

Plural: \emph{die Protokolle}. \emph{Protokoll führen} — to take the minutes.

\begin{cousinweb}
\emph{Protokoll} came through Medieval Latin \emph{prōtocollum} from Byzantine Greek \emph{prōtókollon} (EtymDB), the first sheet glued onto a roll of papers, from \emph{prôtos} \textquotedblleft{}first\textquotedblright{} and \emph{kólla} \textquotedblleft{}glue\textquotedblright{}.
\end{cousinweb}

\begin{grammarlens}[title={What you've built}]
One more word for telling, explaining and describing at B1.
\end{grammarlens}

\subsection*{Your turn}
\begin{itemize}
\raggedright
  \item \emph{Say it:} \emph{das Protokoll}
  \item \emph{Say it:} \emph{das Protokoll}, once more
  \item \emph{Say it:} a sentence of your own with \emph{das Protokoll}
  \item \emph{Recall:} say the German for the meeting, then the German for the session, then say \emph{das Protokoll} again
\end{itemize}

\begin{tcolorbox}[breakable,colback=teal!4,colframe=teal!35!black,title={Before you move on}]
What does das Protokoll mean? (\textquotedblleft{}The minutes, the record; protocol\textquotedblright{}.) Say it once more.
\end{tcolorbox}
\section[präsentieren]{präsentieren — to present}

\label{lesson:GE-C354-praesentieren}



\begin{quote}
Before the new one: say the German for the session, then the German for the minutes.
\end{quote}

\subsection*{You'll want to know: präsentieren}
\textbf{präsentieren} — \textquotedblleft{}to present\textquotedblright{}.

\begin{quote}
Morgen präsentiert sie ihre Ideen vor der ganzen Abteilung.
\end{quote}

\begin{quote}
\emph{Tomorrow she is presenting her ideas to the whole department.}
\end{quote}

Perfekt: \emph{hat präsentiert}, with no \emph{ge-}.

\begin{cousinweb}
\emph{präsentieren} was borrowed in Middle High German as \emph{prēsentieren} from Old French \emph{présenter} (EtymDB), which goes back to Latin \emph{praesentāre} \textquotedblleft{}to show, to present\textquotedblright{}, from \emph{praesēns} \textquotedblleft{}present\textquotedblright{} (Lewis \& Short).
\end{cousinweb}

\begin{grammarlens}[title={What you've built}]
One more word for telling, explaining and describing at B1.
\end{grammarlens}

\subsection*{Your turn}
\begin{itemize}
\raggedright
  \item \emph{Say it:} \emph{präsentieren}
  \item \emph{Say it:} \emph{präsentieren}, once more
  \item \emph{Say it:} a sentence of your own with \emph{präsentieren}
  \item \emph{Recall:} say the German for the session, then the German for the minutes, then say \emph{präsentieren} again
\end{itemize}

\begin{tcolorbox}[breakable,colback=teal!4,colframe=teal!35!black,title={Before you move on}]
What does präsentieren mean? (\textquotedblleft{}To present\textquotedblright{}.) Say it once more.
\end{tcolorbox}
\section[die Präsentation]{die Präsentation — the presentation}

\label{lesson:GE-C354-praesentation}



\begin{quote}
Before the new one: say the German for the minutes, then the German for to present.
\end{quote}

\subsection*{You'll want to know: die Präsentation}
\textbf{die Präsentation} — \textquotedblleft{}the presentation\textquotedblright{}.

\begin{quote}
Ihre Präsentation war kurz und sehr klar.
\end{quote}

\begin{quote}
\emph{Her presentation was short and very clear.}
\end{quote}

Plural: \emph{die Präsentationen}. \emph{Eine Präsentation halten} — to give a presentation.

\begin{cousinweb}
\emph{Präsentation} is a learned borrowing from Late Latin \emph{praesentātiō} \textquotedblleft{}a showing\textquotedblright{}, from \emph{praesentāre} \textquotedblleft{}to show, to present\textquotedblright{}, built on \emph{praesēns} \textquotedblleft{}present\textquotedblright{} (Lewis \& Short).
\end{cousinweb}

\begin{grammarlens}[title={What you've built}]
One more word for telling, explaining and describing at B1.
\end{grammarlens}

\subsection*{Your turn}
\begin{itemize}
\raggedright
  \item \emph{Say it:} \emph{die Präsentation}
  \item \emph{Say it:} \emph{die Präsentation}, once more
  \item \emph{Say it:} a sentence of your own with \emph{die Präsentation}
  \item \emph{Recall:} say the German for the minutes, then the German for to present, then say \emph{die Präsentation} again
\end{itemize}

\begin{tcolorbox}[breakable,colback=teal!4,colframe=teal!35!black,title={Before you move on}]
What does die Präsentation mean? (\textquotedblleft{}The presentation\textquotedblright{}.) Say it once more.
\end{tcolorbox}
